\begin{frame}{Background 1: few-shot prompting}
\begin{block}{Definition}
Few-shot prompting gives the LLM a small number of demonstrations directly in the input prompt. The model is not fine-tuned; it completes the new test instance by following the demonstrated pattern.
\end{block}
\vspace{0.2cm}
\[
  p = \langle x_1 \cdot y_1 \rangle \concat \langle x_2 \cdot y_2 \rangle \concat \dots \concat \langle x_k \cdot y_k \rangle
\]
\[
  \text{Inference input: } p \concat x_{\text{test}} \quad \Rightarrow \quad y_{\text{test}}
\]
\vspace{0.2cm}
\begin{exampleblock}{Intuition}
The prompt teaches the model the expected format and the task behavior from examples instead of changing model parameters.
\end{exampleblock}
\source{PAL paper, Section 2.}
\end{frame}
